\chapter{Describing Wisconsin School Board Elections}
\label{chap:wisco}
\chapterstyle{deposit}






\section{Consider This}

Wisconsin is located in the region of the United States known as the Midwest and ranks 
23rd in terms of total area, 20th in terms of population size with 5.7 million residents, 
and has over 420 local school districts. Recently, Wisconsin has become a crucial swing state 
exhibiting a close split between Republican and Democratic voters in presidential 
politics.\citet{RefWorks:311} 

\section{Background and Legal Basis of School District Elections}

This project represents the first attempt to comprehensively study the patterns 
in school board elections across an entire state over multiple years. Before 
examining the patterns that emerge both within and between school districts, 
it is important to first explore the data available on school board elections. 
Wisconsin, like many other Midwest states and unlike states in the South, is a 
local control state. This means that Wisconsin school districts have revenue 
authority independent of municipalities and counties, and are governed by 
a non-partisan elected board independent of counties and other local government 
bodies. The independence of Wisconsin school districts is established in the 
state constitution. 

This independence gives Wisconsin school districts considerable power and authority. 
Importantly, school boards administer their own elections and retain their own 
election records. This represents a serious barrier to studies of school board 
elections due to the burden associated with collecting and organizing school board 
elections from over 420 administrative entities. This is why previous studies of 
school boards have tended to focus on candidates, random samples of school boards, 
or on states with county-wide school boards such as North Carolina or Georgia. 

Wisconsin school board elections are held consistently on the same dates across 
the state, during the spring election cycle. The spring primary election is the 
third Tuesday in February, and the spring general election is the first Tuesday 
in April (Wis. Stat \S 120.06). Wisconsin school boards must have between 3 and 9 
members with some exceptions made for boards of 11 members (Wis. Stat \S 120.01). 
Board member terms are not consistent across school boards, though 3 year terms are 
the most common with limited exceptions for terms of 1 or 2 years (Wis. Stat \S 120.06 (3)). 
Board members are elected by a plurality and must be residents of the district to 
be eligible for office (Wis. Stat \S 120.06 (2) (a)). All candidate verification, 
ballot preparation, and election announcement is the responsibility of the school district 
clerk (Wis. Stat \S 120.06 (8)). Primary elections occur if there are more than 
twice as many candidates as there are members to be elected (Wis. Stat \S 120.06 (7) (b)). 
School boards may choose to apportion their seats in dramatically different ways 
ranging from a board of 3-9 members each representing the entire district, to a board 
of 3-11 members each representing either a specific region of the school district or 
the entire district. The only restriction on the way seats are allocated is that the 
plan is proposed with a petition by a threshold of voters and passed by the board 
at an annual meeting or through an election (Wis. Stat \S 120.02 (2)). 

School board election records came in many forms. The official and preferred 
records included the Statement of the Board of Canvass and the Certificate of the 
Board of Canvass. These two documents are legally mandated for each school board 
election and certify the winner for each school board seat, as well as the votes 
cast for each candidate. In most cases the records were provided electronically, 
though some were also received as hard copy in the mail. Full details on the data 
collection process can be found in Appendix \ref{app:data}. 
These records were then transcribed from the various records provided by school 
districts into spreadsheets that combined together form the key database for this 
analysis. 

Using the names, vote totals, and winners from each school board election allow 
for the construction of measures of the number of candidates in each election 
cycle, the number of contested seats, the level of voter participation, and other 
metrics of electoral activity at the school board level. Importantly, because the 
data was collected for several years in each district, time-variant trends can be 
evaluated and explored. 


\section{Studying Wisconsin School Board Elections}
\label{sec:wiscolaw}
In Wisconsin, school boards are the offical holder of all school district elections. 
This means that offical school board election records are maintained solely by 
school districts themselves and gathering school board election results in Wisconsin 
requires requesting records from individual school districts.\footnote{See Appendix 
\ref{app:data} and the Appendix \ref{app:log} for details on the request process.} 
An additional challenge is that though school boards are required by law to retain 
these records, the record keeping is often incomplete, unclear, or confusing (Wis. 
Stat \S 120.06 (8) (g)).\footnote{While a few enterprising county clerks have 
collected school board election results and made them available, these can be incomplete 
by only including the precincts that are within the county. They also are not official. 
\citep[personal communication with][]{clerk1}.}In order to conduct this study then, a records 
request was made to all 424 public school districts in the state of Wisconsin. 
This request was e-mailed to the school district administrator in each of Wisconsin's 
424 school districts in January of 2013. Three follow up waves were conducted throughout 
2013, and at the end of the collection period,312 
districts responded with records. The median district provided 
11 
years of records. Overall, records were provided on 4,181 school 
board races involving 5,863 candidates in the state. For each election 
the votes received by each candidate was recorded, as well as whether the candidate 
won, was an incumbent, was a repeat candidate in the school district, or was a 
minor candidate. For each race, whether the candidate was running for district 
wide or regional seat was also recorded. 

While each district was requested to provide the maximum amount of records legally 
required to be retained, ten years, many districts were unable to locate complete 
records going this far back. 

\begin{knitrout}
\definecolor{shadecolor}{rgb}{0.969, 0.969, 0.969}\color{fgcolor}\begin{figure}[]

\includegraphics[width=\maxwidth]{../../phd_figs/chap3-yeargraph} \caption[Number of Districts Providing Records by Year]{Number of Districts Providing Records by Year\label{fig:yeargraph}}
\end{figure}


\end{knitrout}

The good news is that the most crucial period of the sample, from 2007 to 2012, is 
the period with the most complete set of records. 

\begin{knitrout}
\definecolor{shadecolor}{rgb}{0.969, 0.969, 0.969}\color{fgcolor}\begin{figure}[]

\includegraphics[width=\maxwidth]{../../phd_figs/chap3-racesgraph} \caption[Number of Races Provided by Districts]{Number of Races Provided by Districts. Red line indicates the median.\label{fig:racesgraph}}
\end{figure}


\end{knitrout}





The mean district provided records on 13.4 school 
board races. 283 districts had 
a race in each year between 2007 and 2012, while 29 
did not. 

\subsection{General Elections, Primaries, and Special Elections}

General elections for school board members follow the schedule described in Section 
\ref{sec:wiscolaw}. However, for this project, I requested records for \emph{all} elections 
for school board office including school board member primaries and special elections 
if applicable. School board primaries occur whenthe number of candidates exceeds twice 
the number of seats open in a race.\footnote{In Wisconsin first class cities, currently 
only Milwaukee, a primary must be held whenever there are more than two candidates for 
a race. Wis. Stat \S 8.11 (2m)} Special 
elections can also occur if a candidate is recalled under Wis. Stat \S 9.10 (1) (a) 
which states generally that a petition for recall must be signed by electors equal 
in number to at least 25\% of the votes cast for the office of President of the United 
States at the last election within the same jurisdiction.\footnote{It is unclear if this 
would hold for school districts because the determination to hold the election is made by 
school district election official, who may determine that a reliable count of votes cast 
is not available and thus may defer to a formula specified in statute to determine the 
number of petitioners required to initiate a recall. Wis. Stat \S 9.10 (1) (c) to (d). 
Cited in a Staff Brief by the Joint Legislative Council Special Committee on Election 
Law Review, October 2004 \url{http://libcd.law.wisc.edu/~wilc/sb/sb\_2004\_07.pdf}} 
A special election is also required to fill a vacancy for board of school directors 
in the Milwaukee Public School system (MPS). [Wis. Stat \S 8.50 (intro.)] A 
special election is not required to fill a vacancy for any other school district - 
instead vacancies are filled by appointment of the remaining members [Wis. Stat \S 17.26 (1)]. 

For the entire dataset there are 4,000 
general election races,174 
primary election races, and 7
special elections. Table \ref{tab:electype} shows the number of races of each type 
for the period between 2007 and 2012 with the best election record coverage.

% latex table generated in R 3.1.0 by xtable 1.7-3 package
% Sat Jul 26 09:47:37 2014
\begin{table}[ht]
\centering
\begin{tabular}{rrrr}
  \hline
 & General & Primary & Special \\ 
  \hline
2007 & 399 &  15 &   0 \\ 
  2008 & 378 &  13 &   0 \\ 
  2009 & 393 &  15 &   1 \\ 
  2010 & 392 &  16 &   0 \\ 
  2011 & 388 &  12 &   0 \\ 
  2012 & 391 &  22 &   0 \\ 
   \hline
\end{tabular}
\caption{Number of records for general, primary, and special school board elections 2007-2012.} 
\label{tab:electype}
\end{table}


In all, 100 districts 
reported a primary, with 61 
reporting a primary between 2007 and 2012. Only 1
district reported a special election during this period, and 312 
reported a general election race. 

Since Wisconsin school boards are non-partisan offices, a primary is not a symbol of 
an ideological challenge to an incumbent or a battle to win the hearts and minds of 
strict partisans. Instead, a primary is a signal of a substantial level of candidate 
activity for a school board race. 

The only special election that occurred during the 2007 to 2012 period. This was a
recall election held on December 29\^{th}, 2009 in the Crivitz school district which 
was triggered after a school board member, David Kwiatkowski resigned his seat 
that November in the face of a recall petition.\footnote{For a full account of this 
particular case, see the Peshtigo Times article online \url{http://www.peshtigotimes.net/?id=12617}}

\subsection{Contestation: Frequency and Context}

The biggest question in the school board literature seems to be -- are school 
board elections contested, and if so, when? While Chapter \ref{chap:candidacy} 
will explore when candidates emerge for school board elections, this section will 
first look at how often elections are contested, and how this varies compared to 
other elected bodies. 

Unlike studies of congressional or state legislative races, where the barriers to 
appearing on the ballot are generally understood, I begin by discussing the decisions 
a potential school board candidate faces when deciding to run. To appear on the 
ballot in the spring election a candidate may not begin circulating nomination 
papers prior to December 1 in the previous year. The nomination papers must be 
filed no later than 5 p.m. on the first Tuesday in January preceding the election. 
[Wis. Stat \S 8.10 and 8.15] The number of signatures by electors varies in statute, 
depending on the proximity of the school district to cities. In Milwaukee Public Schools 
candidates must obtain between 400 to 800 signatures within the jurisdiction they 
are running to represent (either district wide, or a specific region of the district). 
If the school district contains any territory that lies within a second class city, a 
candidate must obtain between 100 and 200 signatures. School districts that have no 
territory that lies within a first or second class city are free to deterimine 
whether or not require nomination papers.\footnote{Milwaukee is the only first 
class city that meets all requirements in Wis. Stat \S 62.05 (I). There are currently 
sixteen second class cities in Wisconsin, and two more which qualify but have not switched 
status yet. The second class cities are Madison, Green Bay, Kenosha, Racine, Appleton, 
Waukesha, Oshkosh, Eau Claire, Janesville, West Allis, La Crosse, Sheboygan, Wauwatosa, 
Fond du Lac, Brookfield, and Superior.} [Wis. Stat \S 120.06 (a)] The statutes specify that these remaining 
districts can choose to require nomination papers from candidates, but may only require 
between 20 and 100 signatures of electors [Wis. Stat \S 8.10 (3) and 8.15 (6)]. In 
jurisdictions that do not require elector signatures, candidates must file a declaration 
of candidacy to register for the election with the election clerk in the school district. 


Let's consider the pool of school board races for which records were successfully 
obtained. Of the 4,181, 4,000 
were general elections. Another 174 
were primary elections, and the reamining 7 
were special or recall elections. Of the general election races, 
2,855 feature 
two or more candidates who either were on the ballot or garnered more than 20 votes.\footnote{
It is not uncommon in school board races for a write-in candidate to garner more than 20 votes.} 
However, it is important to note that in many school districts voters select more than 
one candidate for each seat and elect the top N vote recipients to the board. A 
better measure of contestation is to see how many races featured more candidates than 
winners - 1,808 
races featured this, or 45.2\% of all races. 

Next, let's look at how often incumbents stand for election and how often they are 
defeated. In order to code incumbency, two data sources were used. First, if a 
candidate appeared in consecutive races and was coded as a winner in the prior 
race, they are coded as an incumbent. However, for races that were the beginning 
of the dataset, incumbency was coded using a master roster of known school board 
members provided by the Wisconsin Department of Public Instruction to the author.\footnote{
This collection is not required or official, but is maintained by the DPI in partnership with 
the Wisconsin Association of School Boards.} This collection indicates name, school district, 
and year of service for school board members in Wisconsin. Candidates appearing on the ballot 
in the early years of a school district's election records were cross referenced against this 
list to infer their incumbency status. 















